\documentclass[11pt,a4paper]{article}
\usepackage[margin=1in]{geometry}
\usepackage{amsmath,amssymb,amsfonts,booktabs,array,tikz}
\usetikzlibrary{shapes,arrows,positioning}

\title{\textbf{Tripartite Arithmetic Architecture}}
\author{Verified Computational Framework}
\date{}

\begin{document}
\maketitle

\section{The Dual-Layer Architecture}

\subsection{The Flux Perimeter ($\mathbf{F}$)}

Residue classes $\{1, 2, 4, 5, 7, 8\}$ modulo 9. Characterized by $\nu_3(n) = 0$ (not divisible by 3).

\textbf{Properties:}
\begin{itemize}
    \item Contains all twin prime pairs $(p, p+2)$ for $p > 3$
    \item Partitioned into three digital-root pairs: $(2,4)$, $(5,7)$, $(8,1)$
    \item Each pair corresponds to $p \pmod{9} \in \{2, 5, 8\}$
\end{itemize}

\textbf{Verified counts} ($p \leq 1{,}000{,}000$):
\begin{center}
\begin{tabular}{cccc}
\toprule
$p \pmod{9}$ & DR Pair & Count & Sum DR \\
\midrule
2 & $(2, 4)$ & 2{,}651 & 6 \\
5 & $(5, 7)$ & 2{,}788 & 3 \\
8 & $(8, 1)$ & 2{,}729 & 9 \\
\bottomrule
\end{tabular}
\end{center}

\textbf{Theorem:} For every twin prime pair $(p, p+2)$ with $p > 3$:
\begin{enumerate}
    \item $p \equiv 2, 5, \text{ or } 8 \pmod{9}$ (Flux Perimeter)
    \item $p + (p+2) = 2(p+1) \equiv 0 \pmod{12}$ (multiple of 12)
    \item $\text{DR}(p + (p+2)) \in \{3, 6, 9\}$ (Invariant Spine)
\end{enumerate}

The mapping $p \pmod{9} \to \text{DR}(\text{sum})$ is deterministic:
\begin{align}
p \equiv 2 \pmod{9} &\implies \text{DR}(p + (p+2)) = 6 \\
p \equiv 5 \pmod{9} &\implies \text{DR}(p + (p+2)) = 3 \\
p \equiv 8 \pmod{9} &\implies \text{DR}(p + (p+2)) = 9
\end{align}

\subsection{The Invariant Spine ($\mathbf{O} \to \mathbf{S}$)}

Residue classes $\{3, 6, 9\}$ modulo 9. Characterized by $\nu_3(n) \geq 1$.

\textbf{Properties:}
\begin{itemize}
    \item Acts as the system's mathematical attractor
    \item All twin prime sums (for $p > 3$) collapse into $\{3, 6, 9\}$
    \item Square transitions: $3^2 = 9 \to \text{DR} = 9$, $6^2 = 36 \to \text{DR} = 9$
    \item Oscillator states $\{3, 6\}$ square-collapse into Singularity $\{9\}$
\end{itemize}

\section{The 10-Sum Identity}

Diametric pairs across the base-10 axis:
\begin{align}
1 + 9 &= 10 \implies \text{DR} = \mathbf{1} \\
2 + 8 &= 10 \implies \text{DR} = \mathbf{1} \\
3 + 7 &= 10 \implies \text{DR} = \mathbf{1} \\
4 + 6 &= 10 \implies \text{DR} = \mathbf{1} \\
5 + 5 &= 10 \implies \text{DR} = \mathbf{1}
\end{align}

All complementary pairs reduce to the unit identity $\text{DR} = 1$.

\section{Reversal Dynamics}

\begin{center}
\begin{tabular}{ccccc}
\toprule
Base & Reversal & Product & DR & State \\
\midrule
37 & 73 & $37 \times 73 = 2701$ & 1 & $\mathbf{F}$ \\
10 & 01 & $10 \times 1 = 10$ & 1 & $\mathbf{F}$ \\
25 & 52 & $25 \times 52 = 1300$ & 4 & $\mathbf{F}$ \\
36 & 63 & $36 \times 63 = 2268$ & 9 & $\mathbf{S}$ \\
\bottomrule
\end{tabular}
\end{center}

\section{Modular Form Locking}

\textbf{E8 Theta Function} ($a(n) = 240\sigma_3(n)$):
\begin{itemize}
    \item Baseline zero rate mod 37: $15.22\%$ ($761/5000$)
    \item When $11 \mid n$: $91.6\%$ zero rate ($416/454$), $6.02\times$ baseline
    \item When $37 \mid n$: $11.1\%$ zero rate ($15/135$)
\end{itemize}

\textbf{Ramanujan Tau Function}:
\begin{itemize}
    \item Baseline zero rate mod 37: $1.76\%$ ($88/5000$)
    \item When $37 \mid n$: $0\%$ zero rate ($0/135$)
    \item When $11 \mid n$: $1.1\%$ zero rate ($5/454$), $0.63\times$ baseline
\end{itemize}

\section{Combinatorial Invariants}

\begin{itemize}
    \item $4! = 24 \implies \text{DR}(24) = 6$
    \item $192 \implies \text{DR}(192) = 1 + 9 + 2 = 12 \implies \text{DR}(12) = 3$
    \item Twin prime count to $10^6$ (standard $6m\pm1$ form): $8{,}168 = 2{,}651 + 2{,}788 + 2{,}729$
    \item Total including special pair $(3,5)$: $8{,}169$
\end{itemize}

\end{document}
